\section{Proceso constructivo}

A continuación se detallarán los procesos constructivos, tanto para el Clarificador Secundario (CS) como para el Campamento de los Trabajadores.

\subsection{Clarificador secundario}

El CS es un estanque circular de hormigón armado con un diámetro interior de 26 metros. Dado que el enunciado no especifica la altura de los muros, se adopta como supuesto de diseño una altura de 4{,}5 m, valor representativo para clarificadores secundarios con tasas de desbordamiento superficial en el rango de 20 a 30 m\textsuperscript{3}/m\textsuperscript{2}/día, habitual en plantas de tratamiento de aguas servidas de mediana y gran escala.

Dada la función hidráulica del CS, tanto su estructura como sus superficies interiores deben ser impermeables. La norma ACI 350-06 \cite{ACI350} establece criterios de fisuración más restrictivos para estructuras de contención de líquidos o residuos ambientales, condición que determina los requerimientos de diseño y proceso constructivo de esta estructura. En base a lo anterior, se adoptan las dimensiones de diseño indicadas en la Tabla~\ref{tab:cs_dim}.

\begin{table}[H]
  \centering
  \caption{Dimensiones de diseño del Clarificador Secundario \textnormal{[SA-01, SA-02]}}
  \label{tab:cs_dim}
  \begin{tabular}{lcc}
    \toprule
    Elemento                 & Dimensión & Unidad \\
    \midrule
    Diámetro interior        & 26{,}0    & m \\
    Diámetro exterior        & 26{,}7    & m \\
    Altura de muros          & 4{,}5     & m \\
    Espesor de muros         & 0{,}35    & m \\
    Espesor de losa de fondo & 0{,}40    & m \\
    \bottomrule
  \end{tabular}
\end{table}

\subsubsection{Moldajes}

Dada la geometría circular y el requerimiento de impermeabilidad del CS, se propone el uso de moldajes trepantes o autodeslizantes (encofrado deslizante) para la ejecución de los muros. Este sistema consiste en un encofrado que avanza de manera continua o semifija en alzada a medida que el hormigón va adquiriendo resistencia suficiente, lo que permite obtener un muro de altura completa con mínimas interrupciones en el proceso de hormigonado.

Las principales ventajas de este sistema para el CS son:

\begin{itemize}
    \item Minimización de juntas de hormigonado: al permitir un vaciado prácticamente continuo, se reduce significativamente la cantidad de juntas frías, mejorando directamente la impermeabilidad de la estructura.
    \item Calidad superficial: se obtiene una superficie interior homogénea, lo que favorece la resistencia química al contacto con aguas servidas.
    \item Optimización de plazos: se eliminan los tiempos muertos asociados al desencofrado y reencofrado por tramos, lo que es especialmente relevante en estructuras de gran altura relativa.
    \item Adaptabilidad a planta circular: los encofrados deslizantes se fabrican en paneles curvos ajustados al radio interior y exterior del muro, garantizando la geometría circular requerida.
\end{itemize}

No obstante, su uso requiere una planificación detallada del proceso de hormigonado, un abastecimiento continuo y sin interrupciones de hormigón fresco, y una supervisión permanente en obra para controlar la velocidad de ascenso del encofrado en función de las condiciones ambientales y la velocidad de fraguado.

Para la losa de fondo, se propone el uso de moldajes tradicionales de madera o metálicos, dado que su geometría es plana y no presenta las exigencias geométricas de los muros. Se construirá en una sola etapa de hormigonado, apoyada directamente sobre una solera de hormigón de limpieza, permitiendo un control adecuado de nivel y planitud.

Cabe destacar que la inversión inicial en encofrado deslizante se justifica con mayor claridad en caso de que el proyecto contemple la construcción de más de un CS, dado que el costo fijo del sistema se amortiza entre las distintas unidades. En cualquier caso, se recomienda realizar un análisis de costo-beneficio comparativo frente a alternativas de encofrado trepante por tramos antes de tomar la decisión final.

\subsubsection{Método y consideraciones de colocación y compactación del hormigón}

\paragraph{Colocación}

Para el vaciado de los muros, se propone el uso de una bomba estacionaria ubicada en el exterior del anillo, con un alcance mínimo de 30 metros, suficiente para cubrir la totalidad del diámetro interior del CS. La tubería de distribución se posicionará de manera que el punto de descarga se desplace perimetralmente conforme avanza el encofrado deslizante, asegurando un llenado uniforme en toda la circunferencia.

Se deben respetar las siguientes consideraciones durante la colocación:

\begin{itemize}
    \item La altura libre de caída del hormigón no debe superar 1{,}5 m, para evitar la segregación de áridos y la pérdida de homogeneidad de la mezcla.
    \item El hormigón debe colocarse en capas horizontales de espesor no mayor a 30--40 cm en el caso de muros, para facilitar la compactación y evitar la formación de bolsones de aire.
    \item El tiempo transcurrido entre la fabricación y la colocación del hormigón debe mantenerse dentro del tiempo de trabajabilidad establecido en la especificación técnica, el cual no deberá exceder los 90 minutos desde la salida de la planta, salvo que se utilicen aditivos retardantes debidamente especificados.
    \item Durante la colocación se debe evitar el desplazamiento horizontal del hormigón con el vibrador, ya que esto puede inducir segregación.
\end{itemize}

Para la losa de fondo, el vaciado se realizará desde el centro hacia la periferia o en franjas paralelas, según lo defina el plan de hormigonado, asegurando continuidad en el frente de avance y evitando juntas frías no planificadas.

\paragraph{Compactación}

Se recomienda la compactación mediante vibradores internos de inmersión (vibradores de aguja), los cuales deben ser operados por personal capacitado siguiendo los parámetros indicados en la Tabla~\ref{tab:vibracion}. En caso de que la densidad de armadura del muro lo requiera, se evaluará la pertinencia del uso de hormigón autocompactante (HAC), el cual no requiere vibración mecánica y facilita el llenado de secciones congestionadas; esta decisión quedará supeditada a la definición del diseño de armadura.

\begin{table}[H]
  \centering
  \caption{Parámetros de vibración interna para muros del Clarificador Secundario}
  \label{tab:vibracion}
  \begin{tabular}{p{6cm} p{8.5cm}}
    \toprule
    \textbf{Parámetro}              & \textbf{Especificación} \\
    \midrule
    Diámetro del vibrador           & 40--50\,mm (compatible con separación entre barras) \\
    Radio efectivo de acción        & 200--300\,mm ($\approx 6 \times$ el diámetro) \\
    Espaciado de inserción          & $\leq 1{,}5 \times$ radio de acción ($\approx 450$\,mm) \\
    Profundidad de penetración      & 100--150\,mm dentro de la capa inferior \\
    Tiempo por punto de inserción   & 5--15\,s \\
    Velocidad de retiro             & $\approx 80$\,mm/s (lenta, para cerrar el canal) \\
    \bottomrule
  \end{tabular}
\end{table}

\subsubsection{Curado del hormigón}

Dado que el CS es una estructura de contención de líquidos, el curado es una etapa crítica para alcanzar la resistencia, impermeabilidad y durabilidad requeridas. Se propone evaluar las condiciones ambientales previo al inicio del hormigonado, con especial atención a la temperatura ambiente, la humedad relativa y la velocidad del viento, variables que en el sector de Salar del Carmen-La Negra (Antofagasta) presentan condiciones habitualmente desfavorables para el curado del hormigón.

En función de las condiciones observadas, se adoptará una de las siguientes estrategias:

\begin{itemize}
    \item Condiciones desfavorables (alta evaporación): se utilizará curado con mantas húmedas o aplicación de compuesto de curado filmógeno conforme a la norma ASTM C309 \cite{ASTMC309}, el cual forma una película impermeable sobre la superficie del hormigón que retiene la humedad interna y previene la evaporación excesiva. Esta solución es especialmente adecuada para la cara superior de la losa de fondo.
    \item Condiciones favorables: se podrá optar por curado húmedo continuo, manteniendo la superficie constantemente humedecida durante un período mínimo de 14 días, en concordancia con los requerimientos de la norma ACI 350-06 para estructuras de contención de líquidos \cite{ACI350}.
\end{itemize}

En ambos casos, el curado debe iniciarse tan pronto como sea posible después del término del vaciado y el acabado superficial, para evitar la contracción plástica del hormigón antes de que este adquiera resistencia.

\subsubsection{Juntas de hormigonado}

\paragraph{Juntas en muros}

Gracias al uso de encofrado deslizante, el hormigonado de los muros puede realizarse de forma continua y sin interrupciones programadas, lo que minimiza la presencia de juntas de hormigonado y mejora la impermeabilidad global de la estructura. No obstante, en caso de que ocurra una interrupción no planificada del proceso (falla mecánica, interrupción del suministro de hormigón, entre otros), se generará una junta fría, la cual deberá ser tratada antes de reanudar el vaciado mediante los siguientes pasos:

\begin{enumerate}
    \item Esperar a que el hormigón colocado alcance resistencia suficiente para no ser perturbado.
    \item Eliminar la lechada superficial y el material suelto mediante chorro de agua a presión o sandblasting.
    \item Humedecer la superficie sin dejar agua estancada (condición SSS, saturado superficialmente seco).
    \item Aplicar un puente de adherencia (lechada de cemento o epoxi, según especificación) antes de reanudar el vaciado.
\end{enumerate}

\paragraph{Junta losa de fondo--muro}

La junta entre la losa de fondo y el muro es la única junta de construcción programada e inevitable en el CS, y corresponde a la zona de mayor solicitación hidrostática de la estructura, por lo que debe recibir un tratamiento especial.

Dado que la presión hidrostática actúa en sentido perpendicular al plano de la junta (es decir, de adentro hacia afuera en la base del muro), se recomienda incorporar una waterstop o banda de estanqueidad de PVC o material hidrófilo en el plano de la junta. Este elemento, embebido en el hormigón de ambas caras de la junta, actúa como barrera física al paso del agua y compensa los movimientos diferenciales entre la losa y el muro.

El procedimiento constructivo para esta junta es el siguiente:

\begin{enumerate}
    \item Hormigonado de la losa de fondo: se ejecuta en primer lugar, dejando una espera de armadura vertical que corresponde al arranque del muro. Durante el vaciado, se posiciona la waterstop centrada en el espesor de la losa, en el plano de encuentro con el futuro muro.
    \item Preparación de la superficie: una vez endurecida la losa, se limpia y rugosifica la superficie de contacto mediante chorro de agua a presión o escarificación mecánica, eliminando la lechada superficial para exponer el árido y mejorar la adherencia.
    \item Puente de adherencia: se aplica una lechada de cemento o mortero epoxi sobre la superficie preparada, en condición SSS, inmediatamente antes de iniciar el vaciado del muro.
    \item Hormigonado del muro: se inicia el vaciado del muro a partir de la junta, asegurando que las primeras capas de hormigón sean compactadas con especial cuidado en la zona de arranque, para evitar coqueras u oclusiones en la zona de mayor presión hidrostática.
\end{enumerate}

Se recomienda que la posición de esta junta se encuentre al menos 15--20 cm por encima del nivel de la losa terminada, de modo que la discontinuidad constructiva no coincida exactamente con el punto de máxima presión en la base del muro.

\subsection{Campamento de los trabajadores}
\label{sec:campamento}

Dado que el enunciado no especifica la cantidad de módulos requeridos, se adopta como supuesto que, al tratarse de un proyecto de gran envergadura con una duración estimada de al menos dos años, la solución más adecuada es un campamento modular prefabricado. Esta alternativa consiste en módulos de estructura metálica con paneles de madera o tipo sándwich, instalados sobre una losa de hormigón, lo que permite una construcción rápida, un buen aislamiento térmico y la posibilidad de desmontaje y reutilización al término del proyecto. Adicionalmente, este tipo de solución resulta especialmente ventajosa en zonas remotas como el sector Salar del Carmen-La Negra, donde la logística de materiales convencionales puede ser compleja y costosa.

\subsubsection{Proceso constructivo}

El proceso constructivo del campamento se puede dividir en las siguientes etapas:

\begin{enumerate}
    \item Preparación del terreno: se realiza el descapote y nivelación del área destinada al campamento, retirando el material orgánico superficial y cualquier elemento que pueda generar asentamientos diferenciales. En caso de suelos blandos o con baja capacidad de soporte, se deberá realizar una mejora de suelo mediante compactación mecánica y/o incorporación de material granular seleccionado.

    \item Subbase granular: sobre el terreno compactado se coloca una capa de subbase granular de al menos 15 cm de espesor, compactada al 95\% del Próctor Modificado, que actúa como capa drenante y de nivelación, asegurando una superficie uniforme para el posterior vaciado de la losa.

    \item Losa de fundación: se construye una losa de hormigón armado que distribuye las cargas de los módulos hacia el suelo. Aunque las cargas involucradas no son de gran magnitud, la losa debe garantizar una superficie nivelada, estable y durable durante toda la vida útil del campamento. Se considera un espesor mínimo de 15 cm, con armadura de temperatura y retracción en ambas direcciones.

    \item Instalaciones embebidas: previo al hormigonado de la losa, se instalan las tuberías de agua potable, alcantarillado y ductos eléctricos que quedarán embebidos o bajo la losa, coordinando su posición con la disposición definitiva de los módulos.

    \item Instalación de módulos: una vez que la losa ha alcanzado la resistencia mínima requerida (típicamente a los 7 días), se procede a la instalación de los módulos prefabricados mediante grúa, los cuales se anclan a la losa mediante pernos de anclaje embebidos durante el hormigonado o mediante fijaciones postinstaladas.

    \item Conexiones de servicio: finalmente, se realizan las conexiones de agua, electricidad, alcantarillado y telecomunicaciones entre módulos y hacia las redes generales del proyecto.
\end{enumerate}

\subsubsection{Colocación y compactación del hormigón}

Para la losa de fundación del campamento se propone el uso de hormigón convencional G25, el cual ofrece una resistencia adecuada para las solicitaciones esperadas y es compatible con las condiciones de exposición del sector (ambiente seco, con posible presencia de sulfatos en el suelo, según lo indicado en la Sección~\ref{sec:especificacion_hormigon}).

El vaciado se realizará mediante camión mixer con canaleta o bomba, dependiendo de la distancia entre el punto de descarga y la zona de vaciado. Dado que se trata de una losa de poco espesor y geometría simple, no se requieren equipos de bombeo de largo alcance. Se recomienda avanzar por franjas paralelas, asegurando continuidad en el frente de vaciado para evitar juntas frías no planificadas.

La compactación se realizará con vibradores internos de aguja de 35 a 50 mm de diámetro, complementados con regla vibrante superficial para garantizar la planitud de la losa y facilitar el acabado superficial. El acabado final se realizará con helicóptero mecánico o llana metálica, dependiendo del nivel de planitud requerido para el apoyo de los módulos.

Una vez finalizado el vaciado, la losa se cubrirá con membrana de curado o lámina de polietileno, que se mantendrá durante al menos 7 días para prevenir la evaporación excesiva. Dado el clima desértico del sector de Antofagasta, con alta radiación solar y vientos frecuentes, se deberá prestar especial atención a esta etapa para evitar la fisuración por retracción plástica antes del fraguado inicial.

\subsubsection{Juntas de contracción}

Dado que la losa de fundación del campamento constituye una gran superficie de hormigón expuesta a condiciones ambientales extremas (alta radiación solar, variaciones térmicas importantes entre el día y la noche propias del clima desértico de Antofagasta), es indispensable controlar la fisuración por retracción mediante la disposición planificada de juntas de contracción.

Las juntas de contracción son planos de debilidad inducidos en la losa que permiten que las fisuras por retracción ocurran de manera controlada y en ubicaciones predefinidas, en lugar de aparecer aleatoriamente en la superficie. Se propone el siguiente esquema:

\begin{itemize}
    \item Espaciado: las juntas se dispondrán en una grilla regular con un espaciado máximo de 4{,}5 m en ambas direcciones, lo que equivale a una relación de aproximadamente 30 veces el espesor de la losa (15 cm), valor recomendado por la norma ACI 360R \cite{ACI360} para losas sobre terreno.
    \item Profundidad: las juntas se ejecutarán con una profundidad mínima de 1/4 del espesor de la losa (aproximadamente 4 cm), suficiente para inducir el plano de debilidad sin comprometer la integridad estructural.
    \item Ejecución: las juntas se realizarán con sierra de disco dentro de las primeras 4 a 12 horas posteriores al vaciado, antes de que se desarrollen tensiones de retracción significativas, pero después de que el hormigón haya endurecido lo suficiente para no desportillarse en los bordes del corte.
    \item Sellado: una vez producida la contracción y estabilizadas las fisuras en las juntas, se procederá al sellado con masilla de poliuretano o silicona flexible, para evitar la infiltración de agua o materiales incompresibles que puedan generar tensiones adicionales.
\end{itemize}

Adicionalmente, se dispondrán juntas de aislación en el perímetro de la losa y en los puntos de encuentro con elementos verticales (muros, columnas o pernos de anclaje), con el fin de desacoplar los movimientos de la losa respecto a los elementos adyacentes.

\subsection{Hormigonado en condiciones extremas}

El sector industrial de Salar del Carmen--La Negra se ubica aproximadamente a 22 km al sureste de la ciudad de Antofagasta, en pleno Desierto de Atacama, a una altitud de aproximadamente 420 m.s.n.m. Las condiciones climáticas de este sector corresponden a un clima desértico árido extremo (clasificación Köppen BWk), caracterizado por una muy baja humedad relativa, precipitaciones prácticamente nulas, alta radiación solar y vientos sostenidos provenientes del suroeste con velocidades habituales de 15 a 30 km/h.

Dado que la construcción de la planta se extenderá por al menos dos años, las faenas de hormigonado se realizarán inevitablemente tanto en período estival como invernal, por lo que es necesario establecer consideraciones específicas para cada condición. Se adoptan los siguientes supuestos de diseño climático para el sector de La Negra:

\begin{table}[H]
  \centering
  \caption{Supuestos de condiciones climáticas extremas para el sector Salar del Carmen--La Negra}
  \label{tab:clima}
  \begin{tabular}{lcc}
    \toprule
    \textbf{Parámetro}               & \textbf{Verano (dic--mar)} & \textbf{Invierno (jun--ago)} \\
    \midrule
    Temperatura máxima diurna        & $\geq 30\,^\circ$C         & $\approx 15$--$18\,^\circ$C \\
    Temperatura mínima nocturna      & $\approx 18\,^\circ$C      & $\leq 5\,^\circ$C           \\
    Humedad relativa                 & $<\,20\%$                  & $<\,30\%$                   \\
    Velocidad del viento             & 15--30 km/h                & 15--30 km/h                 \\
    Radiación solar                  & Alta (cielo despejado)     & Moderada                    \\
    \bottomrule
  \end{tabular}
\end{table}

\noindent Las normas de referencia utilizadas para las consideraciones de hormigonado en condiciones extremas son la NCh170:2016 (Hormigón -- Requisitos generales), el ACI 305R (Guía para hormigonado en clima caluroso) y el ACI 306R (Guía para hormigonado en clima frío).

% -----------------------------------------------------------
\subsubsection{Hormigonado en condiciones de verano (clima caluroso)}

La NCh170:2016 no define el clima caluroso a partir de una temperatura única, sino como una condición de alta evaporación de agua del hormigón, cuyos principales factores son la temperatura ambiente, la humedad relativa del aire, la velocidad del viento y la temperatura del hormigón. De forma complementaria, la guía ACI 305R asocia esta condición a las tasas de evaporación que perjudican la calidad del hormigón fresco, las que en climas normales pueden activarse con temperaturas sobre los 28--30\,$^\circ$C. En el sector de La Negra esta condición se ve agravada permanentemente por la combinación de baja humedad relativa (inferior al 20\%) y vientos sostenidos, por lo que las precauciones de clima caluroso deben considerarse activas durante prácticamente toda la temporada estival, e incluso en días cálidos de primavera u otoño.

Los principales riesgos asociados al hormigonado en clima caluroso son:

\begin{itemize}
    \item Aumento de la demanda de agua de la mezcla, incrementando la relación agua/cemento y reduciendo la resistencia y durabilidad del hormigón.
    \item Aceleración del fraguado, reduciendo el tiempo de trabajabilidad y dificultando la colocación y compactación.
    \item Alta tasa de evaporación superficial, con riesgo de fisuración por retracción plástica antes del fraguado inicial.
    \item Desarrollo de gradientes térmicos al interior de la masa de hormigón, especialmente crítico en elementos de gran espesor como la losa de fondo del CS.
\end{itemize}

\noindent La temperatura máxima admisible del hormigón fresco al momento de la colocación es de 35\,$^\circ$C, conforme a lo establecido por la NCh170:2016 y la guía ACI 305R. Cabe señalar que la NCh170:2016 eliminó el límite de 16\,$^\circ$C que la versión de 1985 imponía a los elementos masivos (dimensión menor superior a 0,80\,m), por lo que la losa de fondo del CS (0,40\,m de espesor) se rige por el límite general de 35\,$^\circ$C. Para cumplir con este límite y mitigar los riesgos descritos, se proponen las siguientes medidas:

\paragraph{Diseño de la mezcla}
\begin{itemize}
    \item Utilizar cementos de bajo calor de hidratación o con adiciones de puzolana/escoria, que reducen el calor generado durante la hidratación y moderan el aumento de temperatura interno.
    \item Incorporar aditivos retardantes de fraguado (conforme a ASTM C494 Tipo B o D) para extender el tiempo de trabajabilidad y compensar la aceleración del fraguado por temperatura, manteniendo la relación agua/cemento especificada en el diseño.
    \item Minimizar el contenido de cemento mediante el uso de adiciones activas (ceniza volante, microsílice), sin comprometer la resistencia ni la impermeabilidad requerida por ACI 350-06 \cite{ACI350}.
\end{itemize}

\paragraph{Producción y transporte}
\begin{itemize}
    \item Enfriar los materiales de la mezcla previo al amasado: utilizar agua fría o con hielo triturado como parte del agua de amasado, y almacenar los áridos bajo sombra o regándolos con agua fría para reducir su temperatura. El agua de mezcla es el componente más eficiente para reducir la temperatura del hormigón fresco, dado su alto calor específico.
    \item Hormigonar en los horarios de menor temperatura diurna, preferentemente en la madrugada o primeras horas de la mañana (entre las 06:00 y las 10:00 hrs), evitando las horas de mayor radiación solar y temperatura ambiente.
    \item Minimizar el tiempo de transporte y espera entre la planta de hormigón y el punto de colocación, de modo que no se supere el tiempo de trabajabilidad especificado ni la temperatura máxima admisible en la descarga. Se recomienda que el tiempo total desde la salida de planta hasta la compactación no exceda los 60 minutos.
    \item Pintar los tambores de los camiones mixer de color blanco o plateado para reducir la absorción de radiación solar durante el transporte.
\end{itemize}

\paragraph{Colocación y compactación}
\begin{itemize}
    \item Humedecer previamente los moldajes, la armadura y cualquier superficie que estará en contacto con el hormigón fresco, para evitar la absorción de agua de la mezcla y la transferencia de calor por conducción.
    \item Proteger las superficies en proceso de vaciado con toldos o cubiertas temporales que impidan la incidencia directa de la radiación solar sobre el hormigón fresco.
    \item Verificar la temperatura del hormigón en la descarga antes de cada vaciado, registrándola en el libro de obra.
\end{itemize}

\paragraph{Curado}
\begin{itemize}
    \item Iniciar el curado tan pronto como sea posible después del término del acabado superficial, sin esperar al fraguado total.
    \item Aplicar compuesto de curado filmógeno (ASTM C309) o cubrir con mantas húmedas y láminas de polietileno para retener la humedad y proteger de la evaporación inducida por viento y radiación.
    \item Mantener el curado durante un mínimo de 14 días para estructuras de contención de líquidos (CS) y 7 días para elementos no expuestos directamente al agua (losa del campamento), conforme a ACI 350 y NCh170 respectivamente.
    \item Monitorear la tasa de evaporación potencial mediante el nomograma de Menzel (o equivalente ACI 305R): si la tasa supera 1,0 kg/m\textsuperscript{2}/h, se deben implementar medidas adicionales de protección superficial, como barreras corta-viento y nebulizadores de agua.
\end{itemize}

% -----------------------------------------------------------
\subsubsection{Hormigonado en condiciones de invierno (clima frío)}

Según la NCh170:2016, se considera que existe una condición de hormigonado en clima frío cuando, durante los tres días previos al hormigonado, se registra una temperatura media diaria menor que 5\,$^\circ$C y la temperatura ambiente es menor o igual que 10\,$^\circ$C por más de 12 horas, continuas o acumuladas, dentro de un período de 24 horas. La guía ACI 306R establece un criterio equivalente. Si bien las temperaturas diurnas en el sector de La Negra durante el invierno oscilan entre los 15 y 18\,$^\circ$C (lo que en principio no constituye un clima frío severo), las temperaturas nocturnas pueden descender a valores cercanos o inferiores a 5\,$^\circ$C, especialmente considerando que el sector se encuentra alejado de la moderación térmica del océano, en un ambiente con cielos despejados y alta irradiación nocturna, propios del desierto de interior.

Por lo tanto, si las faenas de hormigonado se extienden hasta horarios nocturnos o si las temperaturas nocturnas descienden de forma sostenida, se activan las condiciones de hormigonado en clima frío y deben adoptarse las medidas correspondientes. Adicionalmente, la norma NCh170:2016 establece que la temperatura del hormigón al momento de la colocación no debe ser inferior a 5\,$^\circ$C.

Los principales riesgos asociados al hormigonado en clima frío son:

\begin{itemize}
    \item Congelamiento del agua de mezcla antes o durante el fraguado, lo que puede destruir la microestructura del hormigón antes de que este desarrolle resistencia suficiente.
    \item Retardo excesivo del fraguado y endurecimiento, que prolonga los tiempos de desencofrado y puede comprometer el programa de obra.
    \item Desarrollo insuficiente de resistencia a temprana edad, lo que impide la aplicación de cargas en los plazos previstos.
    \item Gradientes térmicos entre el núcleo del hormigón (que genera calor de hidratación) y la superficie expuesta, que pueden inducir fisuras por tensiones térmicas al momento de retirar la protección.
\end{itemize}

\noindent Para mitigar estos riesgos, se proponen las siguientes medidas:

\paragraph{Diseño de la mezcla}
\begin{itemize}
    \item Utilizar cementos de alta resistencia inicial (Tipo III según ASTM C150 o equivalente nacional), que generan mayor calor de hidratación y permiten un desarrollo de resistencia más rápido a bajas temperaturas.
    \item Incorporar aditivos acelerantes de fraguado (ASTM C494 Tipo C o E), que permiten alcanzar la resistencia mínima de protección contra el congelamiento (aproximadamente 3,5 MPa) en un plazo menor.
    \item Reducir la relación agua/cemento al mínimo compatible con la trabajabilidad requerida, lo que también reduce el volumen de agua libre susceptible de congelarse.
\end{itemize}

\paragraph{Producción y transporte}
\begin{itemize}
    \item Calentar el agua de amasado y, de ser necesario, los áridos finos, antes de la mezcla, para garantizar que la temperatura del hormigón fresco al momento de la colocación se encuentre entre 10\,$^\circ$C y 21\,$^\circ$C, rango recomendado por ACI 306R.
    \item No calentar directamente el cemento, ya que puede generar falso fraguado.
    \item Cubrir y aislar los camiones mixer durante el transporte para evitar pérdidas de calor.
    \item Evitar hormigonar sobre superficies congeladas o con hielo: calentar previamente los moldajes, la armadura y el suelo de apoyo hasta alcanzar una temperatura mínima de 0\,$^\circ$C en contacto con el hormigón, conforme a ACI 306R.
\end{itemize}

\paragraph{Protección y curado}
\begin{itemize}
    \item Una vez colocado el hormigón, cubrir inmediatamente las superficies expuestas con mantas de aislación térmica (frazadas de lana de vidrio, mantas de polietileno con cámara de aire u otros sistemas equivalentes) para retener el calor de hidratación y mantener la temperatura del hormigón por encima de 5\,$^\circ$C a 5 cm de la superficie expuesta, conforme al requisito de la NCh170:2016.
    \item En caso de temperaturas muy bajas o vientos fuertes que generen una pérdida de calor excesiva, se deberá implementar un sistema de calefacción activa mediante encierros calefaccionados o calefactores de propano/aire caliente, teniendo precaución de evitar la desecación del hormigón (el calefactor no debe apuntar directamente a la superficie del hormigón).
    \item Mantener la protección durante un período mínimo de 48 a 72 horas hasta que el hormigón alcance una resistencia mínima de 3,5 MPa (resistencia de protección contra el congelamiento), y no retirar la protección abruptamente: la caída de temperatura del hormigón en las 24 horas posteriores al retiro de la protección no debe superar los 5\,$^\circ$C, para evitar fisuras por choque térmico, conforme a ACI 306R.
    \item Continuar el curado húmedo por el período mínimo establecido para cada estructura (7 a 14 días según corresponda), asegurando que el agua de curado no se congele.
\end{itemize}

\paragraph{Monitoreo}
En ambas condiciones extremas (verano e invierno), se debe llevar un registro diario de temperatura del hormigón fresco en la descarga y de temperatura ambiente, humedad relativa y velocidad del viento en obra, para verificar el cumplimiento de las condiciones especificadas y tomar decisiones en tiempo real sobre la necesidad de medidas adicionales de protección.